\section{External Process Supervision and Resource Control}\label{sec:process-isolation}

Media reconstruction combines in-process parsers with supervised external
programs. Static raster images and GIF are decoded and encoded inside the Rust
process. APNG, animated WebP, MP4, WebM, MP3, WAV, FLAC, and Ogg structural
paths use bounded native parsers that retain selected encoded payloads.
Semantic video and audio paths invoke FFprobe and FFmpeg through a process
supervisor. These execution modes have different containment properties and
support different claims.
Structural reconstruction can support container and metadata statements. It
cannot support neutralization of a payload retained inside copied pixel,
sample, or codec data. A policy requiring semantic reconstruction therefore
refuses a structural fallback when a semantic decoder is unavailable.

Every handler dispatch is wrapped by \texttt{catch\_unwind}. A panic becomes a
typed \texttt{HandlerPanic} error and candidate bytes are not returned. Panic
text is reduced to a bounded printable summary before it reaches logs. This
mechanism depends on the checked-in \texttt{panic = "unwind"} release profile.
A consumer that rebuilds the crate with aborting panics loses this containment
property. Unwinding also does not isolate memory corruption in native code,
bound allocator growth automatically, or recover from process termination.
The in-process image and GIF codecs therefore remain in the trusted computing
base for availability and memory safety beyond the guarantees of Rust and
their dependencies.

The external-process supervisor creates temporary input and output files,
passes local paths to the tools, and disables interactive standard input.
FFmpeg receives an explicit protocol allowlist containing only
\texttt{file} and \texttt{pipe}, together with a denylist for network and
indirect protocols. Audio reconstruction selects the first audio stream. Video
reconstruction selects the first video stream and, when present, the first
audio stream. Metadata, chapters, subtitle streams, and data streams are
excluded. The generated file remains a candidate. It is read only
after a successful child exit and then traverses the same output-size,
authorizing-grammar, type-convergence, and signature gates as native output.

\begin{figure*}[t]
\centering
\resizebox{0.96\textwidth}{!}{%
\begin{tikzpicture}[
  lifeline/.style={draw,rounded corners=2pt,minimum width=26mm,minimum height=8mm,font=\small,align=center},
  msg/.style={-{Latex[length=2mm]},thick},
  ret/.style={-{Latex[length=2mm]},dashed},
  note/.style={draw,fill=yellow!9,align=left,font=\scriptsize,inner sep=3pt}
]
\node[lifeline] (caller) at (0,0) {facade};
\node[lifeline] (handler) at (3.5,0) {audio/video handler};
\node[lifeline] (supervisor) at (7.2,0) {process supervisor};
\node[lifeline] (decoder) at (10.8,0) {FFprobe/FFmpeg};
\node[lifeline] (validator) at (14.3,0) {authorizing parser};
\foreach \x in {0,3.5,7.2,10.8,14.3} {\draw[densely dotted] (\x,-0.45) -- (\x,-6.7);}
\draw[msg] (0,-1.0) -- node[above,font=\scriptsize]{dispatch with deadline and token} (3.5,-1.0);
\draw[msg] (3.5,-1.8) -- node[above,font=\scriptsize]{local temporary paths and limits} (7.2,-1.8);
\draw[msg] (7.2,-2.6) -- node[above,font=\scriptsize]{spawn supervised direct child} (10.8,-2.6);
\node[note] at (7.3,-3.45) {poll exit, cancellation, deadline,\\and output-file size};
\draw[ret] (10.8,-4.25) -- node[above,font=\scriptsize]{status and bounded diagnostics} (7.2,-4.25);
\draw[ret] (7.2,-5.0) -- node[above,font=\scriptsize]{candidate bytes or typed error} (3.5,-5.0);
\draw[ret] (3.5,-5.75) -- node[above,font=\scriptsize]{candidate and performed level} (0,-5.75);
\draw[msg] (0,-6.45) -- node[above,font=\scriptsize]{separate complete-output parse} (14.3,-6.45);
\end{tikzpicture}%
}
\caption{Semantic reconstruction sequence. A successful decoder exit creates
only a candidate. Cancellation, timeout, launch failure, oversized output, or
failed validation terminates the path without deliverable bytes.}
\label{fig:decoder-sequence}
\end{figure*}

The supervisor polls the child every 25 milliseconds. During each cycle it
checks whether the monitored output file exceeds the configured byte limit,
whether the shared cancellation token is set, and whether the wall-clock
budget has expired. Standard output and standard error are drained
concurrently to avoid pipe backpressure. Each captured stream is bounded to
16 KiB and later sanitized to remove control characters and bidirectional text
controls. A one-second reader-shutdown bound prevents a surviving helper from
holding a pipe indefinitely. The \texttt{ManagedChild} guard remains armed
until normal completion, so an early return or unwinding path still attempts
termination and reaping.

\begin{table*}[t]
\centering
\caption{Comparison of external-decoder containment by deployment platform.}
\label{tab:platform-containment}
\footnotesize
\begin{tabularx}{\textwidth}{L{0.10\textwidth}L{0.28\textwidth}L{0.28\textwidth}Y}
\toprule
Platform & Enforced by the component & Guarantee not supplied & Deployment condition \\
\midrule
Linux & Process group, \texttt{no\_new\_privs}, disabled core dumps, CPU/file/address-space limits, deadline, cancellation, group kill, and bounded diagnostics & No namespace, seccomp, chroot, container, or mandatory-access-control boundary & Add a dedicated user or stronger platform sandbox when decoder filesystem access is in scope \\
macOS & Temporary-file lifetime, deadline polling, cancellation, output-size monitoring, bounded diagnostics, and direct child termination & No Linux process-group, \texttt{no\_new\_privs}, or resource-limit setup in the compiled branch & Embed in an application sandbox with constrained entitlements and storage access \\
iOS & In-process checks and fail-closed policy handling & Conventional applications cannot spawn the packaged FFmpeg/FFprobe child path & Strict semantic audio/video blocks unless a separately evaluated in-process backend is supplied \\
\bottomrule
\end{tabularx}
\end{table*}

\paragraph{Deployment security contract.}
Four claims remain distinct. Process supervision bounds the direct child's
lifetime, diagnostics, and candidate size during ordinary failure or hang
behavior; on Linux, \texttt{SIGXFSZ} is classified as an exceeded output bound.
Decoder isolation is not supplied: protocol restrictions constrain library
input selection but not arbitrary system calls after code execution is
compromised, so filesystem authority must be narrowed by an enclosing sandbox.
Publication provides non-overwrite and only platform-conditional atomic
visibility. Crash durability is outside the contract because the parent
directory is not synchronized. Thus supervision, isolation, visibility, and
power-loss durability must not be reported as one guarantee.

Cancellation is cooperative for in-process work and supervised for external
work. \texttt{DefenseContext} carries an atomic token, an optional start time,
and a millisecond timeout. Handlers call \texttt{check\_limits} at dispatch and
at format-specific loop boundaries. The facade computes a handler budget that
does not widen a caller's earlier deadline. FFmpeg and FFprobe use the
remaining duration as an upper bound. A cancellation or timeout is a terminal
error, not a warning, and the filesystem publication function is never called
on that result. Cooperative checks cannot preempt a single long-running
in-process codec call, which is another reason not to describe the in-process
path as hard real-time containment.

Temporary files have two roles. Decoder scratch files are named, guard-managed
objects created through the \texttt{tempfile} crate and removed when their
guards drop, subject to host-filesystem behavior. Deliverable publication uses
a separate same-directory temporary file. \texttt{defend\_path} first rejects
input and output paths that identify the same entry and refuses a pre-existing
destination without modifying it. It runs the full byte pipeline, writes only
clean bytes to the temporary file, calls \texttt{sync\_all}, and invokes
\texttt{persist\_noclobber}. The generated suffix must agree with the validated
output profile before staging. The supported claim is non-overwrite: a
pre-existing or concurrently created destination is not replaced. Atomic
visibility depends on the pinned crate version, operating system, and filesystem
primitive and must be verified for the deployment. The parent directory is not
synchronized, so crash durability after power loss is not claimed, as required
by the deployment contract above.

At the result boundary, a blocked candidate may be summarized in an in-memory
quarantine envelope containing its reason, digest, sizes, and stage details.
The envelope is not deliverable, and the current implementation does not persist
quarantined content. Every non-clean verdict returns an empty byte vector, zero
delivered size, and the digest of the empty string. Every exposed interface
enforces the same rule. Parser failure, panic, child launch failure, timeout,
cancellation, policy ambiguity, and authorizing-validation failure therefore
share one no-publication postcondition.
